\section{Theory}
\label{sec:theory}

\subsection{The generalized attenuation law}

Partition the event that two agents give the same answer by whether they are
correct. Two agents can agree by both being right, or by both being wrong in the
same way; a correct and an incorrect agent cannot agree, so the cross term
vanishes exactly. Writing $e_i$ for agent $i$'s marginal error rate and $q_{ij}$
for the probability that two agents who are both wrong chose the same wrong
answer,
\[
  m_{ij} \;=\; (1-e_i)(1-e_j) \;+\; e_i e_j\, q_{ij}.
\]

The binary case is the special instance. With $C=2$ there is exactly one wrong
answer, so $q_{ij} \equiv 1$ identically and the law collapses to
$2m_{ij}-1 = (1-2e_i)(1-2e_j)$, the estimator used in prior work. This is not an
approximation that happens to hold; it is an algebraic identity, and it gives us
a control condition with no free parameters. On BoolQ, where $C=2$, the measured
$q$ is $\numMaeBoolqMc$ in mean absolute error against the multiclass
law---exactness recovered from data, which is the strongest available evidence
that the estimator is implemented as specified.

Applying the binary form when $q < 1$ overstates the quality product by exactly
$e_i e_j (1 - q_{ij})$. The error is second order in the error rates, which is
why it hides on easy benchmarks and surfaces on hard ones: on MATH-500 the binary
law's mean absolute error is $\numMaeMathfiveBinary$ against the multiclass law's
$\numMaeMathfiveMc$. \textbf{[C3, C4]}

The chance level of $q$ for a $C$-ary answer space is $q^{\mathrm{chance}} =
1/(C-1)$, and approximately zero for open-ended answers. This constant is the
reference the whole paper turns on: it is what independent wrong answers look
like, and Section~\ref{sec:limits} shows that an adversary sitting at it is
indistinguishable from independent honest failure.

\subsection{Identifiability and receiver conditioning}

With $q$ known --- binary, at chance, or zero --- the triplet identity recovers
every $e_i$ from $n\geq3$ agents. With $q_{ij}$ free, pairwise agreement gives
$\binom{n}{2}$ equations for $n+\binom{n}{2}$ unknowns: short by exactly $n$ for
every $n$. This is structural unidentifiability, not a sampling problem, and no
amount of data fixes it. Three escapes exist --- a structural assumption on $q$,
the full answer-pattern distribution as Dawid--Skene uses, or ground truth. We
take the first at run time and use the third to check it. \textbf{[C4]}

There is also a gap between the quantity the theory names and the one a gate can
compute, and closing it in the wrong direction is the easiest way to build a
mechanism that validates and then fails. The law is stated in terms of
$P(\text{same}\mid\text{both wrong})$, which needs labels; a receiver can only
compute $P(\text{same}\mid\text{both dissent from me})$. The two coincide only
when the receiver is almost always right, so they are different constants and the
ceilings must be measured on the statistic the gate computes. Ours are
receiver-conditioned throughout: honest coincidence is $\numMmluHonestFloor$ on
four-way multiple choice and $\numGsmkHonestFloor$ on open-numeric answers.
\textbf{[C3]}

\subsection{Coordinatability and the invertibility threshold}

An adversary is exploitable only if it can be \emph{coherent}, which requires a
lie rule its members compute identically from the question alone --- the property
we call \emph{coordinatability}. Which answer spaces permit it, and the direction
in which that expectation failed, is measured in Section~\ref{sec:designlaws}
(Law~3). The threshold follows: invert only above the calibrated honest ceiling
for the answer space. Section~\ref{sec:limits} asks what happens when an
adversary reads that ceiling off the paper and sits underneath it.
